% ===============================
% Worksheet / Manual Template
% ===============================
\documentclass[12pt, a4paper]{article}

% ===============================
% Metadata
% ===============================
\newcommand*{\CourseCode}{MAT221}
\newcommand*{\CourseTitle}{Real Analysis}
\newcommand*{\Chapter}{Infinite Series}
\newcommand*{\Topics}{Convergence of Series}
\newcommand*{\DocDate}{November 16, 2025}

\include{header}

% ==========================================
% Document
% ==========================================
\begin{document}
\FrontPage
\Large

\begin{heading}
	Infinite Series
\end{heading}

\begin{definition}[Infinite Series]
	Let $\{a_n\}_{n=1}^{\infty}$ be a sequence of real numbers.  
	The infinite series generated by this sequence is the formal expression
	\[
	\sum_{n=1}^{\infty} a_n .
	\]
\end{definition}

\clearpage

\begin{heading}
	Convergence of Infinite Series
\end{heading}

\begin{definition}[Convergence of Infinite Series]
	We define the corresponding sequence of partial sums $(s_m)$ by
	\[
	s_m = a_1 + a_2 + a_3 + \cdots + a_m,
	\]
	and say that the series 
	\(
	\sum_{n=1}^{\infty} a_n
	\)
	converges to $S$ if the sequence $(s_m)$ converges to $S$.  
	In this case, we write
	\[
	\sum_{n=1}^{\infty} a_n = S.
	\]
	\vspace{5pt}
	
	If the sequence $\{S_N\}$ does not converge, the series is called divergent.
\end{definition}

\clearpage

\begin{problem}
	Discuss the convergence or otherwise of the series
	\[
	\frac{1}{1\cdot 2} + \frac{1}{2\cdot 3} + \frac{1}{3\cdot 4} + \cdots + \frac{1}{n(n+1)} + \cdots\ \text{to } \infty.
	\]
\end{problem}


\clearpage

\begin{problem}
	Test the convergence or otherwise of the series
	\[
	\sum_{n=0}^{\infty} (-1)^n .
	\]
\end{problem}

\clearpage

\begin{heading}
	Divergence
\end{heading}

\vspace{10pt}

\begin{problem}
	Discuss the convergence or divergence of the harmonic series
	\[
	\sum_{n=1}^{\infty} \frac{1}{n}.
	\]
\end{problem}

\clearpage

\begin{heading}
	Convergence using MCT
\end{heading}

\vspace{10pt}

\begin{problem}
Test the convergence of the series:
\[
\sum_{n=1}^{\infty} \frac{1}{n^{2}}.
\]
\end{problem}

\clearpage

\begin{heading}
	Algebraic Limit Theorem for Series
\end{heading}

\vspace{10pt}

\begin{theorem}[Algebraic Limit Theorem for Series]
Let $\{a_k\}$ and $\{b_k\}$ be sequences of real numbers such that
\[
\sum_{k=1}^{\infty} a_k \quad \text{and} \quad \sum_{k=1}^{\infty} b_k
\]
are both convergent series. Then:

\begin{itemize}
    \item[(i)] The termwise sum series
    \[
    \sum_{k=1}^{\infty} (a_k + b_k)
    \]
    also converges, and its sum is the sum of the individual series:
    \[
    \sum_{k=1}^{\infty} (a_k + b_k)
    = 
    \sum_{k=1}^{\infty} a_k
    + 
    \sum_{k=1}^{\infty} b_k.
    \]

    \item[(ii)] For any constant $c \in \mathbb{R}$ (or $\mathbb{C}$), the scalar multiple series
    \[
    \sum_{k=1}^{\infty} c a_k
    \]
    also converges, and
    \[
    \sum_{k=1}^{\infty} c a_k
    =
    c \sum_{k=1}^{\infty} a_k.
    \]
\end{itemize}
\end{theorem}


\clearpage

\begin{heading}
	Cauchy Criterion for Series
\end{heading}

\vspace{10pt}

\begin{theorem}[Cauchy Criterion for Series]
	The series 
	\[
	\sum_{k=1}^{\infty} a_k
	\]
	converges if and only if, for every $\varepsilon > 0$, there exists an $N \in \mathbb{N}$ such
	\[
	n > m \ge N, \quad \implies \quad \left| a_{m+1} + a_{m+2} + \cdots + a_{n} \right| < \varepsilon.
	\]
\end{theorem}

\clearpage

\begin{theorem}
	If the series 
	\[
	\sum_{k=1}^{\infty} a_k
	\]
	converges, then the sequence $(a_k)$ converges to $0$; that is,
	\[
	a_k \longrightarrow 0 \quad \text{as } k \to \infty.
	\]
\end{theorem}

\clearpage

\begin{heading}
	Convergence of Geometric Series
\end{heading}

\vspace{10pt}

\begin{theorem}
	The geometric series 
	\[
	1 + x + x^{2} + x^{3} + \cdots \ \text{to } \infty
	\]
	satisfies the following:

	\begin{itemize}
		\item it converges if $-1 < x < 1$, i.e.\ if $|x| < 1$;
		\item it diverges if $x \ge 1$;
		\item it oscillates finitely if $x = -1$;
		\item it oscillates infinitely if $x < -1$.
	\end{itemize}
\end{theorem}







\EndPage
\end{document}